% Backup — Leitura detalhada: nenhuma política domina (mensagem-chave já
% incorporada ao slide principal "Fronteira de Pareto CTI-NS")
\begin{frame}[noframenumbering]{[Backup] Nenhuma política domina simultaneamente}
  \begin{center}
    \large Custo, serviço e estabilidade são objetivos \highlight{em conflito}
    -- nenhuma política minimiza CTI e maximiza NS ao mesmo tempo.
  \end{center}
  \vspace{1em}
  \begin{itemize}
    \item \textbf{EOQ} e \textbf{$(s,S)$}: melhor equilíbrio custo-serviço
      agregado (NS\,$=$\,0,94)
    \item \textbf{Jornaleiro}: menor custo bruto, mas NS\,$=$\,0,54
      (abaixo do limiar de viabilidade)
    \item \textbf{SA}: alternativa compacta, bom equilíbrio (NS\,$=$\,0,86)
  \end{itemize}
  \note{
    Resumindo a figura anterior: sobre as 145 séries da Bahia, EOQ
    e (s,S) oferecem o melhor equilíbrio agregado entre custo e serviço,
    ambos com nível de serviço de 0,94. O Jornaleiro tem o menor custo
    bruto, mas seu nível de serviço de 0,54 fica abaixo do limiar mínimo de
    viabilidade de 0,70 -- por isso não é recomendável como política única.
    O SA aparece como alternativa intermediária interessante. Essa leitura
    agregada, porém, esconde diferenças importantes entre perfis de
    demanda, que são exploradas a partir daqui.
  }
\end{frame}
